\section{Versuche}
\subsection{Fadenpendel}
\subsection{Feder-Masse-Pendel}
\subsection{Reflexion, Brechung, Beugung}
\startsubsection{title={Stehende Welle}}
\stopsubsection
\subsection{Polarisation}
% Interferenz am Gitter {{{
\subsection{Interferenz am Gitter/Doppelspalt Licht}

\startbuffer[mp:wellen:versuche:interferenzgitter:aufbau]
\startMPpage

u := 1cm;

pickup pencircle scaled 1pt;

draw (-5u, 0)--(0, 0);
draw (1u, 0)--(4u, 0);
draw (5u, 0)--(7u, 0);

z0=(-5u, 0);

pickup pencircle scaled .5pt;

draw z0--(x0, y0 + 2u);
draw (x0 - .5u, y0 + 1.5u)--(x0 + .5u, y0 + 1.5u);

draw (x0 - .5u, y0 + .5u)--(x0 + .5u, y0 + .5u)--(x0, y0 + 1.5u)--cycle;


z1=(-3.5u, 0);

draw z1{dir 45}..(x1, y1 + 2u);
draw z1{dir 135}..(x1, y1 + 2u);

pickup pencircle scaled 1pt;

z2=(-2u, 0);

draw z2--(x2, y2 + .9u);
draw (x2, y2 + 1.1u)--(x2, y2 + 2u);

pickup pencircle scaled .5pt;

z3=(-.5u, 0);

draw z3{dir 55}..(x3, y3 + 2u);
draw z3{dir 125}..(x3, y3 + 2u);

z4=(2u, 0);

for a=1 upto 10:
drawdot (x4, y4 + (.2u * a)) withpen pensquare scaled 1.5pt;
endfor

z5=(6u, 0);

draw z5--z5 + dir(90) * 2u withpen pencircle scaled 1pt;


label.bot("0cm", z0);
label.bot("3.5cm", z1);
label.top("f50", z1 + (0, 2u));
label.bot("8cm", z2);
label.bot("23cm", z3);
label.top("f100", z3 + (0, 2u));
label.bot("33cm", z4);
label.bot("70cm", z5);

\stopMPpage
\stopbuffer

\placefigure[force][mp:wellen:versuche:interferenzgitter:aufbau]{Aufbau}{
\typesetbuffer[mp:wellen:versuche:interferenzgitter:aufbau]
}

\subsubsection{Beobachtung}

% }}}
% Elektronenbeugungsröhre {{{
\startsubsection{title={Elektronenbeugungsröhre}}
\input notes/physik/aufgaben/2024-10-10-1514-versuchsprotokoll_elektronenbeugungsroehre.tex
\stopsubsection
% }}}
